\documentclass[times,twocolumn,final,authoryear]{elsarticle}
\usepackage{vendor/prletters}
\usepackage{amsmath,amssymb,booktabs,graphicx,url}
\usepackage[hidelinks]{hyperref}
% Draft-only front matter overrides for the licensed 2013 PRL template.
% Body geometry, font size, column widths, caption style remain from prletters.
% Do not print fictitious publisher copyright, receipt dates, or publisher logo.
\makeatletter
\def\ps@PRlettersTitle{\def\@oddhead{\small Pattern Recognition Letters\hfill Research manuscript\quad\thepage}\let\@evenhead\@oddhead\def\@oddfoot{}\let\@evenfoot\@oddfoot}
\renewenvironment{abstract}{\global\setbox\absbox=\vbox\bgroup\hsize=\textwidth\noindent\textbf{Abstract}\par\small\noindent}{\par\egroup}
\def\keyword{\global\setbox\keybox=\vbox\bgroup\hsize=\textwidth\small\noindent\textit{Keywords: }\def\sep{;\space}\ignorespaces}
\def\endkeyword{\par\egroup}
\long\def\MaketitleBox{\resetTitleCounters\begin{flushleft}\Large\bfseries\@title\par\vskip12pt\normalfont\small\elsauthors\par\vskip12pt\unvbox\absbox\vskip8pt\unvbox\keybox\end{flushleft}\vskip12pt}
\makeatother

\journal{Pattern Recognition Letters}
\graphicspath{{figures/}}
\makeatletter
\let\tableinput\@@input
\makeatother
\begin{document}
\begin{frontmatter}
\title{Recoverable but unrealized: the mask-quality gap of prototype-based instance segmentation}
\author{Research manuscript; author information supplied separately}
\begin{abstract}
% CLAIM [B1, A2, C1-C7]: 66.0% of strict argmax-conditioned mask failures are recoverable by a
% coefficient correction inside the frozen prototype space, yet the native GT-free head does not
% realize this. We build a four-layer accounting of where mask quality is lost, quantify the
% recoverable share on the full COCO val2017 population, and exclude seven candidate mechanisms
% one by one: candidate-position selection, supervision quantity, ownership readout, ownership
% through a fixed solver, shared final-layer refitting, post-hoc global correction, oracle-target
% distillation into the native branch, and constrained functional projection. The evidence places
% the bottleneck at the function class / parameterization of coefficient realization rather than at
% information availability, and quantifies the ceiling of correction-style fixes.
A correctly localized instance can still fail a strict mask criterion. We ask whether such mask
failures are recoverable inside the frozen representation and, if so, why the native model does not
recover them. On full COCO val2017 we first account for where mask quality is lost across raw
candidate geometry, class-argmax conditioning, and output retention. We then show that
75.0\% of strict argmax-conditioned mask failures admit a coefficient correction that
reaches Mask75 with the prototype bank frozen, under the project's documented finite-oracle
definition re-derived and independently verified in this study (66.0\% under the original panel's
partially recorded decode conventions; Section~\ref{sec:recoverability}). Seven controlled
comparisons exclude the usual
explanations: oracle-target supervision does not transfer to held-out candidates through the native
coefficient branch, and a constrained functional projection retains less oracle gain than a random
support control. The bottleneck is therefore located in how coefficients are realized, not in
whether the information exists. All readouts are diagnostics on a single frozen checkpoint; no
accuracy claim is made.
\end{abstract}
\begin{keyword}
Instance segmentation \sep Error diagnosis \sep Prototype coefficients \sep Mask quality \sep Frozen-model analysis
\end{keyword}
\end{frontmatter}

\section{Introduction}
% BIND [A1-A4]: AP masks the location of the loss; the three-layer accounting motivates the study.
Modern prototype-based instance segmenters decode masks from a shared prototype bank and per-instance
coefficients \citep{bolya2019yolact}. A candidate can localize its object well and still fail a strict
mask criterion. Standard aggregate accuracy does not say where the quality is lost, and observing
excess foreground does not say which upstream stage produced it. We separate the question into
layers: whether the raw candidate pool contains the capability at all, whether semantic
conditioning keeps it, and what the output stage retains \citep{bolya2020tide}.

Our question is sharper. Among instances that are class-argmax correct and Box75 correct, what share
of mask failures could be repaired by changing only the coefficient vector inside the frozen
prototype space, and if that share is large, why has the native model not realized it? We find
that 75.0\% of the strict population is recoverable in exactly this sense (66.0\% under the original
panel's recorded conventions). The first half is a recoverability fact; the second half is an
exclusion problem. We answer both, and the answers change where a fix has to act.

The paper's claim is deliberately negative in a useful way: the recoverable share is large, and
none of the seven candidate mechanisms we test accounts for the gap. Short-budget supervision
toward oracle directions does not transfer to held-out candidates; a constrained functional
projection of the same oracle directions protects reliable regions extremely well but retains only
about a third of the available gain, and less than a random-support control. No test yields a
deployable method; each is a controlled comparison that removes one explanation.

\section{Related work}
% BIND: literature boundaries must be stated as in the project evidence map; no novelty claims on
% structure (prototype-coefficient readout, boundary refinement, error analysis all have precedent).
Prototype--coefficient decoding is established \citep{bolya2019yolact}. Error-analysis tooling for
detection and segmentation is established \citep{bolya2020tide}. Boundary refinement and quality
scoring are established \citep{kirillov2020pointrend,tang2021bpr,huang2019maskscoring,ke2022transfiner}.
Ownership and competition between instances has been modeled before \citep{ke2021bcnet}.
We therefore claim no structural novelty. What is new here is the combination: a strict
argmax-conditioned failure population, a frozen-prototype recoverability measurement on the full
validation set, and a chain of controlled exclusions that together localize the gap.

\section{Accounting: where mask quality is lost}
\subsection{Four evaluation units}
% BIND [A1, A2, A3]: raw pool, argmax conditioning, output retention; identities kept separate.
Let a raw candidate be identified by image, branch, pyramid level, grid position and input shape.
For each ordinary GT we ask, in order: (i) does any raw position achieve Box75 and Mask75; (ii) does
any raw whose predicted-class argmax equals the GT class do so; (iii) does the normal output
retain a good candidate. The units answer different questions and are never interchanged. This
yields the population used below: instances that are argmax-correct and Box75-correct yet have no
argmax-correct Mask75 candidate anywhere.

\subsection{Population and baseline}
% BIND [A1, A2, A4]: 36,335 GT; 6,459 strict mask failures (17.78%); baseline AP.
On COCO val2017 \citep{lin2014coco} with 5{,}000 images and 36{,}335 ordinary GT, the strict
population contains 6{,}459 instances (17.78\%; A1 gives the score-independent geometry
classes, A4 the baseline AP). All readouts use the official frozen YOLO26m-Seg checkpoint
\citep{jocher2026yolo26}. Output-stage retention accounts for a further, separately recorded loss of
otherwise good masks (A3). Table~\ref{tab:accounting} collects the three layers. The remainder of
the paper studies only the strict population.

\begin{table*}[t]
\caption{Where mask quality is lost, in three separate accounting layers (COCO val2017, 5{,}000
images, 36{,}335 ordinary GT, frozen YOLO26m-Seg). The layers answer different questions and are
never interchanged. Shares are of 36{,}335 GT except in the output-retention block, where they are
of the 13{,}568 GT whose normal output fails Mask75.}
\label{tab:accounting}
\centering
\begin{tabular}{@{}llrr@{}}
\toprule
layer & state & GT & share \\
\midrule
raw pool (score-free) & same raw Box75+Mask75        & 23{,}805 & 65.52\% \\
                      & good box, no good mask in pool& 6{,}688 & 18.41\% \\
                      & good mask, no good box       & 487 & 1.34\% \\
                      & good box/mask on different raws& 531 & 1.46\% \\
                      & neither good                  & 4{,}824 & 13.28\% \\
\midrule
argmax-conditioned    & semantic fail                 & 44 & 0.12\% \\
                     & box fail                      & 6{,}335 & 17.43\% \\
                     & \textbf{mask fail (strict)}   & \textbf{6{,}459} & \textbf{17.78\%} \\
                     & box/mask misaligned           & 484 & 1.33\% \\
                     & success                       & 23{,}013 & 63.34\% \\
\midrule
output retention      & no raw Mask75 (of 13{,}568 output failures) & 11{,}512 & 84.85\% \\
                     & good mask, class score $\le$0.001 & 1{,}797 & 13.24\% \\
                     & not in head top-300            & 153 & 1.13\% \\
                     & per-category maxDets budget   & 103 & 0.76\% \\
                     & COCO matching competition     & 3 & 0.02\% \\
\midrule
baseline              & Mask AP / Mask75 matched     & \multicolumn{2}{r}{43.518 / 22{,}767 (62.66\%)} \\
\bottomrule
\end{tabular}
\end{table*}

\section{Recoverability inside the frozen prototype space}
\subsection{Oracle definition}
% BIND [B1 + disclosure]: finite oracle under fixed P; definition and鍙ｅ緞 per the frozen record.
Fix the prototype bank, the coefficient of the candidate, its predicted box and the decoding path.
For each instance solve, in the GT-box region of the input grid,
\begin{equation}
\Delta c^{*}_{\lambda}=\arg\min_{\Delta c}\
\frac{1}{A}\sum_{u}\operatorname{BCE}\!\big(p_u^{\mathsf T}(c_0+\Delta c),y_u\big)
+\frac{\lambda}{2}\|\Delta c\|_2^2,\qquad
c^{*}=c_0+\Delta c^{*}_{\lambda},
\label{eq:oracle}
\end{equation}
with $\lambda$ fixed in advance; Eq.~\ref{eq:oracle} is the project's recorded finite-oracle
convention and is GT-assisted by construction. An instance is \emph{recoverable} when some
argmax-correct raw's $c^{*}$ reaches Mask75 after the official decode.

\subsection{Result}
\label{sec:recoverability}
% BIND [B1, B2, B4, B3]: re-derived 75.02% (4,843/6,456) under the documented 7D definition;
% recorded 66.00% (4,263/6,459) under the original panel's partially recorded decode conventions.
We re-derived the whole panel under Eq.~\ref{eq:oracle} with the letterbox geometry, GT raster and
decode path rebuilt from the retained per-image records, and verified the re-derivation by
independent re-solves from different initializations, exact re-aggregation, and identity
cross-checks. The re-derived share is \textbf{4{,}843 of 6{,}456 instances (75.0\%)};
Table~\ref{tab:recoverability} reports strata by GT size next to the original panel's recorded
values. The original panel's solver is not part of the retained record, and its decode and ROI
conventions are only partially documented; we therefore report both values and use neither as the
sole basis of any claim. Three GTs (0.05\%) whose GT box degenerates on the decode grid are
recorded as skipped. The difference concentrates in small instances and is consistent with a
threshold-boundary convention effect near the 0.75 criterion, not with a different oracle family.

An independent panel re-measures the same construct. The remainder of the non-recoverable set
splits into fixed-prototype/decoding-limited and support-infeasible instances in the original record
(2{,}169 and 27 respectively), and 309 candidate rows of that record store the support ceiling
rather than a solved value; these accounting facts do not change either recoverable count. Position
selection is not the explanation: across instances with several argmax-correct raws, the oracle
best exceeds the per-instance median by only 1.293 points on average (1.854 on the 4{,}504
instances that have multiple raws). Whatever convention is used, the recoverable share is a large
majority of the strict population, and no fitted mechanism below realizes it.

\begin{table}[b]
\caption{Recoverability of strict argmax-conditioned mask failures inside the frozen prototype
space: an oracle coefficient correction (Eq.~\ref{eq:oracle}) reaches Mask75. For each stratum:
the definition-complete re-derivation of this study (primary) and the original panel's recorded
values (partially recorded decode conventions; both reported per the frozen protocol). Three
degenerate GTs are excluded from the re-derived population.}
\label{tab:recoverability}
\centering
\begin{tabular}{@{}lrrrr@{}}
\toprule
& \multicolumn{2}{c}{re-derived (this study)} & \multicolumn{2}{c}{recorded (original panel)} \\
\cmidrule(lr){2-3}\cmidrule(lr){4-5}
stratum & recov. & share & recov. & share \\
\midrule
small ($<$32$^2$)     & 2{,}817 & 72.5\% & 2{,}285 & 58.8\% \\
medium (32$^2$--96$^2$) & 1{,}266 & 72.7\% & 1{,}237 & 71.0\% \\
large ($\ge$96$^2$)   & 760 & 91.7\% & 741 & 89.4\% \\
\midrule
all                   & \textbf{4{,}843} & \textbf{75.0\%} & 4{,}263 & 66.0\% \\
\bottomrule
\end{tabular}
\end{table}

\section{Why the native model does not realize it}
\label{sec:exclusions}
% BIND [C1-C7]: each exclusion states what was tested, the number, and the scope limit.
Seven comparisons each remove one explanation. Table~\ref{tab:exclusions} summarizes; we state each
with its boundary.

\emph{Supervision quantity.} Adding a spatial-logit teacher toward oracle logit changes moves the
training objective but not the output: the controlled comparison is +0.013 points of image-macro IoU
against its same-budget control, with the target group's interval crossing zero and no additional
Mask75 crossings.

\emph{Ownership readout and solver.} Reading pixel ownership from frozen neck features does not beat
a same-capacity control; routing external spatial evidence through the fixed solver stays below
the native model.

\emph{Shared final layer.} Re-fitting the shared affine final layer on frozen features recovers
46.8\% of the diagnostic objective gap and, when frozen and applied to normal images, yields a
bootstrap-positive real-mask improvement (+0.055 image-macro IoU), concentrated in P4/P5. The
remaining per-instance oracle gap is not recoverable by any shared map.

\emph{Post-hoc global correction.} A frozen global corrector improves image-macro IoU by 0.092
points, CI $[+0.022,+0.167]$, with no reliable Mask75 change.

\emph{Oracle-target distillation.} Fine-tuning only the coefficient heads toward oracle-derived
targets does not transfer: on held-out candidates both the target arm and a same-budget BCE-only
control move \emph{against} the oracle direction (cosine $-0.03$; $-0.12$ at P5 and for large
instances), the repair gain is identical with and without the oracle target, and both damage
previously correct instances at about 7\%.

\emph{Constrained functional projection.} Projecting oracle corrections onto a trust-region
subspace (evidence pixels by current-mask uncertainty, protected pixels by margin) suppresses
protected-zone logit change to 6\% of the oracle's own, but retains 30\% of the oracle IoU gain
against an 80\% gate, and retains less than a random same-cardinality support control (59\%).

\emph{Reading.} The gap is not information availability (the recoverable share is large) and not
any of the seven tested mechanisms. It sits at the realization stage: the native coefficient head's
function class, as trained, does not select the recoverable directions.

\begin{table*}[t]
\caption{Exclusion chain: each row removes one candidate explanation of why the native model does
not realize the recoverable coefficient directions. Numbers are image-macro or per-instance
readouts with bootstrap intervals reported in the project record; scope limits are part of each
result.}
\label{tab:exclusions}
\centering
\begin{tabular}{@{}llp{44mm}p{32mm}@{}}
\toprule
layer & verdict & key number & scope limit \\
\midrule
position selection   & excluded & oracle best--median 1.293 points (1.854 on multi-raw instances) & diagnostic, fixed panel \\
supervision amount  & excluded & teacher vs same-budget control $+0.013$ pp, target-group CI crosses 0 & 253 dev images, 1{,}816 candidates \\
ownership readout   & excluded & true-ROI vs equal-capacity control: IoU $-0.001$, CI 0 & 8$\times$8 soft ownership, 196 images \\
ownership via solver& excluded & external evidence through fixed solver $-4.99$ pp macro IoU & fixed solver parameters, one external model \\
shared final layer  & partial   & 46.8\% of diagnostic gap; $+0.055$ image-macro IoU, CI positive & mask-BCE diagnostic on frozen features \\
global correction   & bounded   & $+0.092$ pp, CI $[+0.022,+0.167]$ & 504 images, 3{,}560 fixed positives \\
oracle distillation & excluded & held-out alignment $-0.03$ (P5 $-0.12$); repair gain equals control; damage 6.7--7.3\% & 1 seed, 3 epochs, 796 fit / 197 dev images \\
constrained projection & excluded & retention 30.2\% vs random support 58.8\% (gate 80\%) & fixed margins and penalty grid \\
\bottomrule
\end{tabular}
\end{table*}

\section{What the gap implies for corrections}
% BIND [A5, A6, A7]: upper bounds and object selection, used to bound the value of correction-style fixes.
Two independent bounding results frame what a correction can buy. Removing false-positive pixels
from already-correctly-boxed predictions raises mask AP by 10.999 points in a GT-assisted ideal
repair, and the failure rate varies strongly with the GT's fill ratio (62.4\% vs 10.5\%
between the lowest and highest fill strata), which is an object-selection signal, not a
mechanism. Decode-level substitutions do not recover the gap. Together with
Section~\ref{sec:exclusions} (Table~\ref{tab:exclusions}) this localizes the value of any
correction-style approach: it can remove pixels, but it cannot supply the coefficient directions
the native head fails to realize.

\section{Limitations}
% BIND: the frozen record's disclosures; nothing beyond them.
Single checkpoint and single architecture family; COCO val2017 images were used repeatedly during
development, so no result is a blind test; oracle quantities are GT-assisted diagnostics and not
method performance; candidate--GT associations are diagnostic, not COCO matching; no accuracy is
claimed. The recoverability measurement's exact objective and the projection audit's margins are
frozen in the project record and cited accordingly. Cross-model generality is untested.

\section{Conclusion}
Mask-quality failures in a prototype-based segmenter are largely recoverable inside its own frozen
representation, and the native prediction does not realize them. The exclusion chain in this paper
moves the question from ``is there information'' to ``how must coefficients be realized'', and
gives the next method attempt a specific target and a set of controlled comparisons it must beat.

\section*{{Data and code availability}}
All readouts use the frozen official YOLO26m-Seg checkpoint and the public COCO validation split
with original annotations; no weights are trained in this study. Every number in
Tables~\ref{tab:accounting}--\ref{tab:exclusions} traces to a retained run record in the project's
experiment archive. Per experiment we retain the protocol frozen before execution, per-candidate
and per-GT records, independent verification runs with their checks, source snapshots with hashes,
and transfer manifests; a machine-readable claim-to-record map (evidence map) enumerates each
claim, its source records, and its verification status. The recoverability re-derivation and its
verification are retained as separate runs of the same archive.

\bibliographystyle{elsarticle-harv}
\bibliography{references}
\end{document}

